% !PRINTED 339
\weberpage{339}
%<<CONT
\gesperrt{$\alpha$, then one obtains a narrower interval $(\alpha', \beta')$
with the same properties, if one}
\weq{4}{
\alpha' = \alpha - \frac{f(\alpha)(\beta - \alpha)}{f(\beta) - f(\alpha)}, \quad \beta' = \beta - \frac{f(\beta)}{f'(\beta)}
}
\gesperrt{sets.}

To prove these results of geometric intuition analytically,
we assume that in the interval $f'(x)$ and
$f''(x)$ are positive, hence $\alpha < \beta$, and
\[
f(\alpha) < 0, \quad f(\beta) > 0.
\]

We restrict ourselves to the consideration of this one
case, since the three others are to be treated in exactly the same
way.

We form the function
\[
\varphi(x) = \frac{f(\beta) - f(x)}{\beta - x}
\]
and its first derivative
\[
\varphi'(x) = \frac{-(\beta - x)f'(x) + f(\beta) - f(x)}{(\beta - x)^2}.
\]
The numerator of this expression
\[
\psi(x) = -(\beta - x)f'(x) + f(\beta) - f(x)
\]
vanishes for $x = \beta$, and its derivative is
\[
\psi'(x) = -(\beta - x)f''(x),
\]
hence negative in the interval. From this it follows that $\psi(x)$ decreases
in the interval as $x$ increases, and therefore, because it vanishes for the
greatest value $x = \beta$, remains positive. Consequently
$\varphi'(x)$ is also positive, and $\varphi(x)$ increases steadily in the interval
as $x$ increases. We accordingly have
\weq{5}{
\begin{gathered}
\frac{f(\beta) - f(\alpha)}{\beta - \alpha} < \frac{f(\beta) - f(x)}{\beta - x} < f'(\beta) \\
\alpha < x < \beta.
\end{gathered}
}
If we now set
\[
\alpha' = \alpha - \frac{f(\alpha)(\beta - \alpha)}{f(\beta) - f(\alpha)}, \quad \beta' = \beta - \frac{f(\beta)}{f'(\beta)},
\]
then
\weq{6}{
\beta' - \alpha' = f(\beta)\left\{\frac{\beta - \alpha}{f(\beta) - f(\alpha)} - \frac{1}{f'(\beta)}\right\},
}
and from (5) it follows that this difference is positive, hence
\[
\alpha < \alpha' < \beta' < \beta.
\]
% 22*
